\section{Util --- Odometer}
%%% generated by the blueprint pipeline from TCSlib.Programs.Util.Odometer
\subsection*{Util --- Odometer: Introduction}

Declarations of the Lean module \texttt{TCSlib.Programs.Util.Odometer}.

\subsection*{Util --- Odometer: Results}

\begin{definition}[Incrementing a digit list]
\lean{Light.incDigits}
\label{Light.incDigits}
\leanok
For a natural number base $b$, define the increment of a list of natural-number digits
recursively: the empty list remains empty; for a first digit $d$ followed by a list $l$,
replace $d$ by $d+1$ and keep the rest when $d+1<b$, and otherwise replace it by zero
and increment the remaining list.
\end{definition}

\begin{definition}[Base expansion digits]
\lean{Light.digitsLE}
\label{Light.digitsLE}
\leanok
For natural numbers $b,n,c$, define the length-$n$ list whose entry at index $k<n$ is
the remainder of $c/b^k$ upon division by $b$.
\end{definition}

\begin{theorem}[Length of the digit list]
\lean{Light.length_digitsLE}
\label{Light.length_digitsLE}
\leanok
\uses{Light.digitsLE}
\difficulty{1}
For all natural numbers $b,n,c$, the list of the first $n$ base-$b$ digits of $c$, with
the digit at position $k$ given by $\lfloor c/b^k\rfloor\bmod b$, has length $n$.
\end{theorem}

\begin{theorem}[Bound on base digits]
\lean{Light.lt_of_mem_digitsLE}
\label{Light.lt_of_mem_digitsLE}
\leanok
\uses{Light.digitsLE}
\difficulty{3}
For natural numbers $b,n,c$ with $b>0$, every entry of the list whose $k$th term, for
$0\le k<n$, is the remainder of $c$ upon division by $b^k$ and then by $b$ is less than
$b$.
\end{theorem}

\begin{theorem}[Digits after increasing the length]
\lean{Light.digitsLE_succ}
\label{Light.digitsLE_succ}
\leanok
\uses{Light.digitsLE}
\difficulty{3}
For natural numbers $b,n,c$, the list of the first $n+1$ base-$b$ remainders of $c$ is
the list whose first term is $c\bmod b$ and whose remaining terms are the first $n$
remainders of $\lfloor c/b\rfloor$.
\end{theorem}

\begin{theorem}[Incrementing a base expansion]
\lean{Light.incDigits_digitsLE}
\label{Light.incDigits_digitsLE}
\leanok
\uses{Light.digitsLE, Light.digitsLE_succ, Light.incDigits, Nat.succ_div_mod_of_eq, Nat.succ_div_mod_of_lt}
\difficulty{5}
For a positive integer base $b$, let the length-$n$ digit list of $c$ have entries
$\lfloor c/b^k\rfloor \bmod b$ for $k=0,\ldots,n-1$, and increment a digit list by
adding one to its first digit when that remains below $b$, or otherwise setting it to
zero and carrying to the rest of the list. Then incrementing the length-$n$ digit list
of $c$ in this way gives the length-$n$ digit list of $c+1$.
\end{theorem}
